% Official CC BY4.0 published HTML S4; original math alttext and author text, literal original numbering.


\par\medskip\noindent\textbf{4 Representations of Lie algebras}\par




In this section, we provide an introduction to Lie algebras, universal enveloping algebras, and their representation theory, following mainly [32, 42].
Afterwards, we introduce and explain the results needed to construct the projectors in Theorem 1.1. These results are typically stated in the language of abstract Lie algebras, and therefore we take time to build up this general theory. In our case of interest, the Lie algebra arises from $\textnormal{GL}_{k}$ and many parts simplify.
We will regularly indicate this in the text for clarity and concreteness.
We establish homogeneous polynomials and metamonomials as a Lie algebra representation, which we will use as running examples.





A Lie algebra $\mathfrak{g}$ is a vector space endowed with a bilinear bracket operation $[-,-]$ which is skew-commutative, i.e., $[X,Y]=-[Y,X]$, and satisfies a relation called the Jacobi identity $[X,[Y,Z]]+[Z,[X,Y]]+[Y,[Z,X]]=0$; in particular, the bracket operation is usually not associative. Every associative algebra $\mathcal{A}$ is a Lie algebra with the commutator bracket given by $[a,b]\coloneqq ab-ba$ for every $a,b\in\mathcal{A}$. In this paper, we are interested in the representation theory of the general linear group $\textnormal{GL}_{k}$, and we study it via the representation theory of its Lie algebra $\mathfrak{gl}_{k}$: this is the space of $k\times k$ matrices endowed with the commutator bracket.





A Lie algebra representation is a vector space $V$ with an associated linear map $\rho:\mathfrak{g}\to\textnormal{End}(V)$ such that $\rho([X,Y])=[\rho(X),\rho(Y)]$ where the bracket on the right is the commutator bracket. We equivalently say that $\mathfrak{g}$ (or its elements) acts on $V$. For $v\in V$, we often write $X.v\coloneqq\rho(X)v$ as a shorthand.
A simple example is $V=\mathbb{C}^{k}$, where $\mathfrak{gl}_{k}$ acts by matrix-vector multiplication. This is called the standard representation of $\mathfrak{gl}_{k}$.
In this paper we examine the space of metapolynomials $\mathbb{C}[\mathbb{C}[x_{1},\ldots,x_{k}]_{d}]_{\delta}$ as a Lie algebra representation of $\mathfrak{gl}_{k}$. The Lie algebra action is explicitly given in in ˜4.2.





Representation theory of Lie algebras is used to study the representation theory of Lie groups because it allows one to “linearize” a group action.
More precisely, every Lie group $G$ has an associated Lie algebra $\mathfrak{g}$ and every Lie group representation (a non-linear map) defines naturally a Lie algebra representation (a linear map).
All Lie algebra representations we will encounter arise from this correspondence. We will describe the $\mathfrak{gl}_{k}$ action on polynomials and metapolynomials explicitly in ˜4.1 and ˜4.2, which may be taken as the definition. The details of the correspondence between Lie groups and Lie algebras are not crucial for the rest of the paper. We briefly outline them in Remark 4.3.





For a vector space $V$ with basis $v_{1},\ldots,v_{r}$, let $v_{1}^{*},\ldots,v_{r}^{*}\in V^{*}$ be the dual basis of $V^{*}$.
Explicitly, $v_{\ell}^{*}$ is the linear form mapping a vector $v=\sum_{i}c_{i}v_{i}\in V$, with $c_{i}\in\mathbb{C}$, to the coefficient $c_{\ell}$.
The basis dual to the standard basis of $\mathbb{C}^{k}$ is the set of variables $x_{1},\ldots,x_{k}$.
Let $E_{i,j}\in\mathfrak{gl}_{k}$ be the $k\times k$ matrix having $1$ at the entry $(i,j)$ and zero elsewhere. Regarding $E_{i,j}:\mathbb{C}^{k}\to\mathbb{C}^{k}$ as a linear map, we have $x_{\ell}\circ E_{i,j}=0$ if $\ell\neq i$ and $x_{\ell}\circ E_{i,j}=x_{j}$ if $\ell=i$; here $\circ$ denotes the composition of functions. Since $\{E_{i,j}\}_{i,j}$ forms a basis for $\mathfrak{gl}_{k}$, the action of $\mathfrak{gl}_{k}$ on a representation $V$ is uniquely determined by the action of the $E_{i,j}$’s.





\par\medskip\noindent\textbf{$\vartriangleright$ Claim 4.1 (Action on polynomials).}\par



The vector space $\mathbb{C}[x_{1},\ldots,x_{k}]_{d}$ of homogeneous degree $d$ polynomials is a representation of $\mathfrak{gl}_{k}$. It is the linearization of the action of $\textnormal{GL}_{k}$ acting by base change on the variables $x_{1},\ldots,x_{k}$.
For $f\in\mathbb{C}[x_{1},\ldots,x_{k}]_{d}$, the action is defined by the shifted partial derivative



\[\begin{aligned}
\displaystyle E_{i,j}.f=-\sum_{\ell=1}^{k}(x_{\ell}\circ E_{i,j})\frac{\partial}{\partial x_{\ell}}f=-x_{j}\frac{\partial}{\partial x_{i}}f.\qquad\text{(1)}
\end{aligned}\]


In particular, for a monomial $f=x_{\ell_{1}}\cdots x_{\ell_{d}}$, we have



\[\begin{aligned}
E_{i,j}.f=\textstyle\sum_{p=1}^{d}x_{\ell_{1}}\cdots x_{\ell_{p-1}}(E_{i,j}.x_{\ell_{p}})x_{\ell_{p+1}}\cdots x_{\ell_{d}}.
\end{aligned}\]





For a multiindex $\mu\in\mathbb{N}^{k},\sum_{i}\mu_{i}=d$, let $c_{\mu}$ be the corresponding variable, that is, the linear map sending a polynomial $f$ to the coefficient of $x^{\mu}$. Let $e_{1},\ldots,e_{k}\in\mathbb{N}^{k}$ be the standard basis vectors: $e_{i}$ is the vector with $1$ at the $i$-th entry and $0$ elsewhere.





\par\medskip\noindent\textbf{$\vartriangleright$ Claim 4.2 (Action on metapolynomials).}\par



The vector space of metapolynomials of format $(\delta,d,k)$, that is $\mathbb{C}[\mathbb{C}[x_{1},\ldots,x_{k}]_{d}]_{\delta}=\mathbb{C}[c_{\mu}:\mu\in\mathbb{N}^{k},|\mu|=d]_{\delta}$ is a representation of $\mathfrak{gl}_{k}$. It is the linearization of the action of $\textnormal{GL}_{k}$ acting by base change on the variables $x_{1},\ldots,x_{k}$.
For a metapolynomial $\Delta\in\mathbb{C}[c_{\mu}:\mu\in\mathbb{N}^{k},|\mu|=d]_{\delta}$, the action is given by





\[\begin{aligned}
\displaystyle E_{i,j}.\Delta=\sum_{\mu}(E_{i,j}.c_{\mu})\frac{\partial}{\partial c_{\mu}}\Delta,\qquad\text{(2)}
\end{aligned}\]


where the summation is over all multi-indices $\mu\in\mathbb{N}^{k}$ with $\sum_{i}\mu_{i}=d$. The action on metavariables is given by



\[\begin{aligned}
E_{i,j}.c_{\mu}=\begin{cases}(\mu_{i}+1)\,c_{\mu+e_{i}-e_{j}}&\text{ if }i\neq j,\\ \mu_{i}\,c_{\mu}&\text{ otherwise},\end{cases}\qquad\text{(3)}
\end{aligned}\]


where we set $c_{\mu+e_{i}-e_{j}}=0$ if $\mu_{j}=0$.
In particular, for a metamonomial $\Delta=c_{\mu_{1}}\cdots c_{\mu_{\delta}}$, we have



\[\begin{aligned}
E_{i,j}.\Delta=\textstyle\sum_{p=1}^{\delta}c_{\mu_{1}}\cdots c_{\mu_{p-1}}(E_{i,j}.c_{\mu_{p}})c_{\mu_{p+1}}\cdots c_{\mu_{\delta}}.
\end{aligned}\]





The following remark illustrates the correspondence between representations of complex algebraic Lie groups and representations of the corresponding Lie algebras. As mentioned before, it is not essential for the rest of the paper, but the proof of ˜4.1 and ˜4.2 are built on this correspondence.





\par\medskip\noindent\textbf{$\blacktriangleright$ Remark 4.3.}\par



A complex algebraic Lie group is a group with the structure of a complex algebraic variety, [32, Ch. 7]. One can show that $\textnormal{GL}_{k}$ is such a group. Let $V$ be a finite-dimensional vector space. A representation of a complex algebraic Lie group $G$ on $V$ is a group homomorphism $\rho:G\to\textnormal{GL}(V)$ which is a regular map, in the sense that the entries of $\rho(g)\in\textnormal{GL}(V)$ regarded as a matrix in a fixed basis, are polynomial functions on $G$.





As a vector space, the Lie algebra of $G$ is, by definition, the tangent space of $G$ at its identity element $\mathfrak{g}=T_{\mathsf{id}}G$. The bracket operation is defined via the adjoint representation of $G$ on $\mathfrak{g}$, see [32, Ch. 8] In the case of $\textnormal{GL}_{k}$ one has $\mathfrak{gl}_{k}\simeq T_{\mathsf{id}_{k}}\textnormal{GL}_{k}$ is the space of $k\times k$ matrices, and the bracket is the standard commutator.





Every representation $\rho:G\to\textnormal{GL}(V)$ of $G$ defines a representation $\mathfrak{g}$ via its differential at the identity: $\textnormal{d}\rho:T_{\mathsf{id}}G\to T_{\mathsf{id}}\textnormal{GL}(V)$; this is a linear map and defines a Lie algebra homomorphism from $T_{\mathsf{id}}G=\mathfrak{g}$ to $T_{\mathsf{id}}\textnormal{GL}(V)=\textnormal{End}(V)=\mathfrak{gl}(V)$.





Explicitly, the value of $\textnormal{d}\rho$ at $X\in\mathfrak{g}$ is the element $\textnormal{d}\rho(X)\in\mathfrak{gl}(V)\simeq\textnormal{End}(V)$ defined by



\[\begin{aligned}
\big{(}\textnormal{d}\rho(X)\big{)}v=\frac{\textnormal{d}}{\textnormal{d}s}\Big{|}_{s=0}\Big{(}\rho\big{(}\gamma(s)\big{)}v\Big{)}\quad\text{ for every $v\in V$};\qquad\text{(4)}
\end{aligned}\]


here $\gamma:\mathbb{C}\to G$ is any smooth curve such that $\gamma(0)=\mathsf{id}_{G}$ and $\gamma^{\prime}(0)=X$. It is easy to see that $\textnormal{d}\rho(X)$ does not depend on the choice of $\gamma$.






\par\medskip\noindent\textbf{Proof of ˜4.1 and ˜4.2.}\par




We derive the Lie algebra representations from their respective group representations, as in Remark 4.3.





Every representation $\rho:G\to\textnormal{GL}(V)$ defines a pullback representation on the space of polynomials of degree $p$ on $V$, $\mathbb{C}[V]_{p}$. For $g\in G$ and $f\in\mathbb{C}[V]_{p}$, this is given by $g\cdot f=f\circ\rho(g^{-1})$ where $\rho(g^{-1}):V\to V$ is regarded as a linear map and $\circ$ denotes the composition of functions.





The standard representation of $\textnormal{GL}_{k}$ on $\mathbb{C}^{k}$ induces in this way a representation on the space of polynomials $\mathbb{C}[x_{1},\ldots,x_{k}]_{d}$. Similarly, the action of $\textnormal{GL}(\mathbb{C}[x_{1},\ldots,x_{k}]_{d})$ induces a representation on the space of metapolynomials $\mathbb{C}[\mathbb{C}[x_{1},\ldots,x_{k}]_{d}]_{\delta}$; the action of $\textnormal{GL}_{k}$ on metapolynomials is the composition of the two representations



\[\begin{aligned}
\textnormal{GL}_{k}\to\textnormal{GL}(\mathbb{C}[x_{1},\ldots,x_{k}]_{d})\to\textnormal{GL}(\mathbb{C}[\mathbb{C}[x_{1},\ldots,x_{k}]_{d}]_{\delta}).
\end{aligned}\]


Now, if $\rho:G\to\textnormal{GL}(V)$ induces the Lie algebra representation $\textnormal{d}\rho:\mathfrak{g}\to\textnormal{GL}(V)$, the pullback representation of $G$ on $\mathbb{C}[V]_{p}$ induces the Lie algebra representation described as follows. For $X\in\mathfrak{g}$, let $\gamma(s)$ be a smooth curve in $G$ such that $\gamma(0)=\mathsf{id}_{G}$ and $\gamma^{\prime}(0)=X$. The chain rule of derivatives implies $\frac{\textnormal{d}}{\textnormal{d}s}\big{|}_{s=0}\gamma(s)^{-1}=-X$. We show this in the case where $G$ is a group of matrices: in this case, one can multiply elements of $G$ by elements of $\mathfrak{g}$ and obtain



\[\begin{aligned}
\frac{\textnormal{d}}{\textnormal{d}s}\bigg{|}_{s=0}\gamma(s)^{-1}=-\bigl{[}\frac{\textnormal{d}}{\textnormal{d}s}\bigg{|}_{s=0}\gamma(s)\bigr{]}\cdot\bigl{[}(\gamma(s)^{-2})\big{|}_{s=0}\bigr{]}=-\gamma^{\prime}(0)\cdot\gamma(0)^{-2}=-X\cdot\mathsf{id}^{-2}=-X.
\end{aligned}\]


In general, the multiplication of matrices should be replaced by the action of $G$ on $\mathfrak{g}$.





Let $V^{*}$ be the dual space of $V$ with basis $z_{1},\ldots,z_{r}$; in particular $\mathbb{C}[V]_{p}$ is the space of polynomials of degree $p$ in $z_{1},\ldots,z_{r}$. By Equation 4, for every $F\in\mathbb{C}[V]_{p}$, $X.F$ is the polynomial function described as follows, on every $\zeta\in V$:



\[\begin{aligned}
\displaystyle X.F(\zeta)= & \displaystyle\frac{\textnormal{d}}{\textnormal{d}s}\Big{|}_{0}(\gamma(s)\cdot F)(\zeta)= \\
 & \displaystyle\frac{\textnormal{d}}{\textnormal{d}s}\Big{|}_{0}(F\circ\gamma(s)^{-1})(\zeta)=\qquad(\text{using chain rule}) \\
 & \displaystyle\sum_{\ell=1}^{r}\frac{\partial}{\partial z_{\ell}}F(\zeta)\cdot(z_{\ell}(\frac{\textnormal{d}}{\textnormal{d}s}\Big{|}_{0}\gamma(s)^{-1}\zeta))= \\
 & \displaystyle\sum_{\ell=1}^{r}\frac{\partial}{\partial z_{\ell}}F(\zeta)\cdot(z_{\ell}(-X.\zeta))=-\sum_{\ell=1}^{r}\frac{\partial}{\partial z_{\ell}}F(\zeta)\cdot z_{\ell}(X.\zeta).
\end{aligned}\]


Specialize the calculation above to the case of polynomials and metapolynomials to obtain Equation 1 and Equation 2. It remains to show Equation 3, which follows from a direct calculation: regarding $c_{\mu}$ as a function of a polynomial $f$, we have



\[\begin{aligned}
E_{ij}.c_{\mu}(f)=-c_{\mu}(E_{ij}.f)=c_{\mu}(x_{j}\frac{\partial}{\partial x_{i}}f)=\begin{cases}(\mu_{i}+1)\,c_{\mu+e_{i}-e_{j}}(f)&\text{ if }i\neq j,\\ \mu_{i}\,c_{\mu}(f)&\text{ otherwise}.\end{cases}\ 
\end{aligned}\]


$\hfill\blacktriangleleft$





Let $V$ be a representation of a Lie algebra $\mathfrak{g}$. A subspace $V^{\prime}\subseteq V$ that is itself a $\mathfrak{g}$-representation is called a subrepresentation. A representation $V$ is irreducible if it contains no subrepresentations except the zero space of $V$ and $V$ itself. A map $\phi:V\to W$ between two representations is called equivariant if it commutes with the action of $\mathfrak{g}$. Schur’s Lemma characterizes equivariant maps between irreducible representations:





\par\medskip\noindent\textbf{Lemma 4.4 (Schur’s Lemma, [32, Lem. 1.7]).}\par



Let $V,W$ be irreducible representations for a Lie algebra $\mathfrak{g}$ and let $\phi:V\to W$ be an equivariant map. Then either $\phi=0$ or $\phi$ is an isomorphism. Moreover, if $V=W$, then $\phi=c\cdot\mathsf{id}_{V}$ for some $c\in\mathbb{C}$.






A complex algebraic Lie group is called reductive if every representation is completely reducible, namely it splits into direct sum of irreducible subrepresentations; the general linear group $\textnormal{GL}_{k}$ is reductive, see e.g. [45, Ch.X]. In the following $\mathfrak{g}$ is the Lie algebra of a reductive group, and explicitly we are interested in the case where $G=\textnormal{GL}_{k}$ and $\mathfrak{g}=\mathfrak{gl}_{k}$.






\par\medskip\noindent\textbf{4.1 Cartan subalgebras and weights}\par




A subalgebra $\mathfrak{t}\subseteq\mathfrak{g}$ is called abelian if $[\mathfrak{t},\mathfrak{t}]=0$. A Cartan subalgebra $\mathfrak{h}$ of $\mathfrak{g}$ is a maximal abelian subalgebra with the property that if $[X,\mathfrak{h}]\subseteq\mathfrak{h}$ then $X\in\mathfrak{h}$. The subalgebra of diagonal matrices in $\mathfrak{gl}_{k}$ is a Cartan subalgebra of $\mathfrak{gl}_{k}$. It turns out that if $G$ is reductive, then all Cartan subalgebras of $\mathfrak{g}$ have the same dimension; in fact they are all equivalent under a natural action of $G$ on $\mathfrak{g}$ [32, Sec. D3]. In the case of $\textnormal{GL}_{k}$, the action on $\mathfrak{gl}_{k}$ is given by conjugation: $g\cdot X=gXg^{-1}$ for $g\in\textnormal{GL}_{k}$ and $X\in\mathfrak{gl}_{k}$. In this case, one can show that any Cartan subalgebra is equivalent to the subspace of diagonal matrices.





If $\mathfrak{g}$ is the Lie algebra of a complex reductive Lie group $G$, then Cartan subalgebras of $\mathfrak{g}$ correspond to maximal abelian subgroups in $G$. For a fixed Cartan subalgebra $\mathfrak{h}$, let $\mathbb{T}$ be the corresponding abelian subgroup of $G$. If $\mathfrak{h}$ is the subset of diagonal matrices in $\mathfrak{gl}_{k}$, then $\mathfrak{h}$ is the Lie algebra of the abelian group $\mathbb{T}$ of diagonal matrices with nonzero diagonal entries.
Irreducible $\mathbb{T}$-representations are $1$-dimensional and they are in one-to-one correspondence with a lattice $\Lambda$ in the space $\mathfrak{h}^{*}$, called the weight lattice of $\mathfrak{g}$; the elements of $\Lambda$ are called weights. When $\mathbb{T}$ is the subgroup of $\textnormal{GL}_{k}$ of diagonal matrices, then this definition coincides with the one used in Section 1 for metapolynomials; see also Example 4.7. The construction of the weight lattice is not straightforward for arbitrary Lie algebras, but it is very explicit in the case where $\mathfrak{h}$ is the subalgebra of diagonal elements in $\mathfrak{gl}_{k}$:





\par\medskip\noindent\textbf{$\vartriangleright$ Claim 4.5.}\par



The weight lattice of $\mathfrak{gl}_{k}$ with respect to the Cartan subalgebra $\mathfrak{h}$ of diagonal elements is the abelian group of integer linear combinations of the diagonal elements, that is, the group of linear functions



\[\begin{aligned}
\displaystyle\alpha:\mathfrak{h} & \displaystyle\to\mathbb{C} \\
\displaystyle\left(\begin{smallmatrix}h_{1}&&\\ &\ddots&\\ &&h_{k}\end{smallmatrix}\right) & \displaystyle\mapsto\alpha_{1}h_{1}+\cdots+\alpha_{k}h_{k}.
\end{aligned}\]


with $\alpha_{1},\ldots,\alpha_{k}\in\mathbb{Z}$. In the following, we usually identify the weight $\alpha$ with the $k$-tuple of its integer coefficients $\alpha=(\alpha_{1},\ldots,\alpha_{k})$.






\par\medskip\noindent\textbf{Proof.}\par




Let $\mathbb{T}$ be the subgroup of diagonal matrices in $\textnormal{GL}_{k}$. Let $(\alpha_{1},\ldots,\alpha_{k})\in\mathbb{Z}^{k}$ be an integer vector.



\[\begin{aligned}
\displaystyle\rho:\mathbb{T} & \displaystyle\to\mathbb{C}^{\times}\simeq\textnormal{GL}_{1} \\
\displaystyle\left(\begin{smallmatrix}t_{1}&&\\ &\ddots&\\ &&t_{k}\end{smallmatrix}\right) & \displaystyle\mapsto t_{1}^{\alpha_{1}}\cdots t_{k}^{\alpha_{k}}
\end{aligned}\]


is a group homomorphism, hence it defines a $1$-dimensional (irreducible) representation of $\mathbb{T}$. It is not hard to prove that all irreducible representations of $\mathbb{T}$ arise in this way.





The induced $\mathfrak{h}$-representation is



\[\begin{aligned}
\displaystyle\textnormal{d}\rho:\mathfrak{h} & \displaystyle\to\mathbb{C}\simeq\mathfrak{gl}_{1} \\
\displaystyle\left(\begin{smallmatrix}h_{1}&&\\ &\ddots&\\ &&h_{k}\end{smallmatrix}\right) & \displaystyle\mapsto\alpha_{1}h_{1}+\cdots+\alpha_{k}h_{k}.
\end{aligned}\]


In other words, irreducible $\mathfrak{h}$-representations induced from $\mathbb{T}$-representations are in one-to-one correspondence with functions $\alpha\in\mathfrak{h}^{*}$ mapping a diagonal element $H\in\mathfrak{h}$ to an integer linear combination of its diagonal entries. $\hfill\vartriangleleft$





Every $\mathfrak{g}$-representation $V$ restricts to an $\mathfrak{h}$-representation, hence it decomposes into the direct sum of irreducible $1$-dimensional $\mathfrak{h}$-representations. Since they are uniquely determined by their weights, we have $V=\bigoplus_{\alpha\in\Lambda}V_{\alpha}^{\oplus m_{\alpha}}$; in particular $m_{\alpha}\neq 0$ only for finitely many $\alpha\in\Lambda$. If $m_{\alpha}\neq 0$, then $\alpha$ is called a weight of $V$, $m_{\alpha}$ is called multiplicity and the direct summand $V_{\alpha}^{\oplus m_{\alpha}}$ is called the weight subspace of $V$ of weight $\alpha$. An element $X\in\mathfrak{h}$ acts simultaneously on the $m_{\alpha}$ copies of $V_{\alpha}$ via scalar multiplication by $\alpha(X)\in\mathbb{C}$. The elements of a weight subspace are called weight vectors; a basis of $V$ consisting of weight vectors is called a weight basis.





\par\medskip\noindent\textbf{Example 4.6.}\par



Consider the $\mathfrak{gl}_{3}$-representation $V=\mathbb{C}[x_{1},x_{2},x_{3}]_{2}$ with the action described in ˜4.1. Let $\mathfrak{h}\subseteq\mathfrak{gl}_{3}$ be the subalgebra of diagonal matrices.
Consider the six monomials of $V$, which define a basis:



\[\begin{aligned}
\begin{array}[]{lll}x_{1}^{2},&x_{2}^{2},&x_{3}^{2},\\ x_{2}x_{3},&x_{1}x_{3},&x_{1}x_{2}.\end{array}
\end{aligned}\]


One can verify that they form a weight basis of $V$. For instance, given $H=\left(\begin{smallmatrix}t_{1}&&\\ &t_{2}&\\ &&t_{3}\end{smallmatrix}\right)=t_{1}E_{11}+t_{2}E_{22}+t_{3}E_{33}\in\mathfrak{h}$, we have



\[\begin{aligned}
H.(x_{1}^{2})=\textstyle\sum_{i}t_{i}E_{ii}.(x_{1}^{2})=-\sum_{i}t_{i}\sum_{j}x_{j}\displaystyle\frac{\partial}{\partial x_{j}}(x_{1}^{2})=-(2t_{1})x_{1}^{2}
\end{aligned}\]


showing that $x_{1}^{2}$ is a weight vector of weight $(-2,0,0)$. More generally, we have



\[\begin{aligned}
\begin{array}[]{rlrlrl}H.(x_{1}^{2})=\hskip-4.62497pt&\hskip-4.62497pt-2t_{1}\,x_{1}^{2},&H.(x_{2}^{2})=\hskip-4.62497pt&\hskip-4.62497pt-2t_{2}\,x_{2}^{2},&H.(x_{3}^{2})=\hskip-4.62497pt&\hskip-4.62497pt-2t_{3}\,x_{3}^{2},\\ H.(x_{2}x_{3})=\hskip-4.62497pt&\hskip-4.62497pt-(t_{2}+t_{3})\,x_{2}x_{3},&H.(x_{1}x_{3})=\hskip-4.62497pt&\hskip-4.62497pt-(t_{1}+t_{3})\,x_{1}x_{3},&H.(x_{1}x_{2})=\hskip-4.62497pt&\hskip-4.62497pt-(t_{1}+t_{2})\,x_{1}x_{2}.\end{array}
\end{aligned}\]


This shows that $V$ decomposes into the sum of six weight spaces, each spanned by one of the weight vectors above. The corresponding weights are the six elements of $\mathfrak{h}^{*}$ mapping $H$ to the eigenvalues of $H$ on each weight vector. As $k$-tuple of integers, they are



\[\begin{aligned}
\begin{array}[]{lll}(-2,0,0),&(0,-2,0),&(0,0,-2),\\ (0,-1,-1),&(-1,0,-1),&(-1,-1,0).\end{array}
\end{aligned}\]


More generally, in a space of polynomials $V=\mathbb{C}[x_{1},\ldots,x_{k}]_{d}$, the monomials form a weight basis. The weight of a monomial $x^{\mu}=x_{1}^{\mu_{1}}\cdots x_{k}^{\mu_{k}}$, regarded as a vector of integers, is $-\mu$.






\par\medskip\noindent\textbf{Example 4.7.}\par



Consider the $\mathfrak{gl}_{2}$-representation $V=\mathbb{C}[\mathbb{C}[x_{1},x_{2}]_{3}]_{2}$ of metapolynomials of format $(2,3,2)$, with the action described in ˜4.2. One can verify that the metamonomials form a weight basis. For instance, given $H=\left(\begin{smallmatrix}t_{1}&\\ &t_{2}\end{smallmatrix}\right)=t_{1}E_{11}+t_{2}E_{22}\in\mathfrak{h}$, we have



\[\begin{aligned}
H.(c_{30}^{2})=\textstyle\sum_{i=1}^{2}t_{i}E_{ii}.(c_{30}^{2})=\textstyle\sum_{i=1}^{2}t_{i}\sum_{\mu}E_{ii}.c_{\mu}\cdot\displaystyle\frac{\partial}{\partial c_{\mu}}(c_{30}^{2})=3\cdot 2\cdot t_{1}\cdot c_{30}^{2}=(6t_{1})c_{30}^{2}.
\end{aligned}\]


showing that $c_{30}^{2}$ is a weight vector of weight $(6,0)$. Similarly, one can compute the weight of every metamonomial:



\[\begin{aligned}
\begin{array}[]{rlrlrl}H.(c_{21}^{2})=\hskip-7.39998pt&(4t_{1}+2t_{2})c_{21}^{2},&H.(c_{12}^{2})=\hskip-7.39998pt&(2t_{1}+4t_{2})c_{12}^{2},&H.(c_{03}^{2})=\hskip-7.39998pt&(6t_{2})c_{03}^{2},\\ H.(c_{30}c_{21})=\hskip-7.39998pt&(5t_{1}+\phantom{1}t_{2})c_{30}c_{21},&H.(c_{30}c_{12})=\hskip-7.39998pt&(4t_{1}+2t_{2})c_{30}c_{12},&H.(c_{30}c_{03})=\hskip-7.39998pt&(3t_{1}+3t_{2})c_{30}c_{03},\\ H.(c_{21}c_{12})=\hskip-7.39998pt&(3t_{1}+3t_{2})c_{21}c_{12},&H.(c_{21}c_{03})=\hskip-7.39998pt&(2t_{1}+4t_{2})c_{21}c_{03},&H.(c_{12}c_{03})=\hskip-7.39998pt&(\phantom{1}t_{1}+5t_{2})c_{21}c_{12}.\end{array}
\end{aligned}\]


Therefore $V$ has seven weight spaces of weight $(p,6-p)$ with $p=0,\ldots,6$. The weight spaces of weight $(4,2),(3,3),(2,4)$ have dimension $2$, the others have dimension $1$.





More generally, the metamonomials give a weight basis for the space of metapolynomials of format $(\delta,d,k)$. A metamonomial $c_{\mu_{1}}\cdots c_{\mu_{\delta}}$ has weight $\mu_{1}+\cdots+\mu_{\delta}$,
in accordance with what we defined in the introduction.








\par\medskip\noindent\textbf{4.2 Roots and highest weights}\par




An important representation is given by the action of the Lie algebra on itself, called the adjoint representation $\textnormal{ad}:\mathfrak{g}\to\textnormal{End}(\mathfrak{g})$ and defined by $\textnormal{ad}(X)=[X,-]$.
By definition, the weight space of weight $0$ is $\mathfrak{h}$ itself; moreover, one can prove that all other weight spaces are $1$-dimensional [32, Ch. 14]. The nonzero weights of the adjoint representation are called the roots of the Lie algebra $\mathfrak{g}$ (with respect to $\mathfrak{h}$); let $\Phi_{\mathfrak{g}}\subseteq\Lambda$ denote the set of roots of $\mathfrak{g}$. It turns out that $\eta\in\mathfrak{h}^{*}$ is a root if and only if $-\eta$ is a root and $\pm\eta$ are the only integer multiples of $\eta$ that are roots.





We compute the roots for $\mathfrak{gl}_{k}$ with $\mathfrak{h}$ the algebra of diagonal matrices. Let $H=\textnormal{diag}(t_{1},\ldots,t_{k})\in\mathfrak{h}$. Then



\[\begin{aligned}
\displaystyle\textnormal{ad}(H)E_{i,j}=HE_{i,j}-E_{i,j}H=t_{i}E_{i,j}-t_{j}E_{i,j}=(t_{i}-t_{j})E_{i,j}.\qquad\text{(5)}
\end{aligned}\]


So $\mathbb{C}E_{i,j}$ is the weight space corresponding to the weight $(e_{i}-e_{j})\in\mathfrak{h}^{*}$. Since the matrices $E_{i,j}$ span $\mathfrak{gl}_{k}$, we obtain that the set of roots is $\Phi_{\mathfrak{gl}_{k}}=\{e_{i}-e_{j}\mid i,j\in[k],i\neq j\}$.





The roots of the Lie algebra shift the weights of the Lie algebra representations: more precisely, let $V$ be a representation of $\mathfrak{g}$, $v\in V$ a weight vector of weight $\lambda$ and $X\in\mathfrak{g}_{\alpha}$ for some root $\alpha$; then $X.v$ is a weight vector of weight $\lambda+\alpha$.





A choice of positive roots is a subset $\Phi_{\mathfrak{g}}^{+}$ of $\Phi_{\mathfrak{g}}$ with the following two properties: for every root $\alpha$, either $\alpha$ or $-\alpha$ belong to $\Phi_{\mathfrak{g}}^{+}$; for every pair of roots $\alpha_{1},\alpha_{2}$ such that $\alpha_{1}+\alpha_{2}$ is a root, if $\alpha_{1},\alpha_{2}\in\Phi_{\mathfrak{g}}^{+}$ then $\alpha_{1}+\alpha_{2}\in\Phi_{\mathfrak{g}}^{+}$. For $\mathfrak{gl}_{k}$, one such choice is $\Phi_{\mathfrak{gl}_{k}}^{+}=\{e_{i}-e_{j}\mid i<j\}$.





All choices of positive roots are equivalent under the action of a finite group $\mathcal{W}_{\mathfrak{g}}$, called the Weyl group of $\mathfrak{g}$. This group acts on $\mathfrak{h}$ and hence on $\mathfrak{h}^{*}$; the action of $\mathfrak{h}^{*}$ sends the set of roots to itself. All possible sets of positive roots are obtained one from the other via the action of $\mathcal{W}_{\mathfrak{g}}$. The Weyl group of $\mathfrak{gl}_{k}$ is the permutation group $\mathfrak{S}_{k}$ on $k$ elements. If $\mathfrak{h}$ is the Cartan subalgebra of diagonal matrices, then $\mathcal{W}_{\mathfrak{g}}$ acts by permuting the diagonal elements; the induced action on $\mathfrak{h}^{*}$ permutes the weights $e_{i}$: for $\sigma\in\mathfrak{S}_{k}$, $\sigma\cdot e_{i}$ is the weight $e_{\sigma^{-1}(i)}$.





Since for every root $\alpha$, either $\alpha$ or $-\alpha$ is positive, the weight decomposition of $\mathfrak{g}$ gives $\mathfrak{g}=\mathfrak{h}\oplus\bigoplus_{\alpha\in\Phi_{\mathfrak{g}}^{+}}(\mathfrak{g}_{\alpha}\oplus\mathfrak{g}_{-\alpha})$. In the case of $\mathfrak{gl}_{k}$, this decomposition is $\mathfrak{gl}_{k}=\mathfrak{h}\oplus\bigoplus_{i<j}(\mathbb{C}E_{ij}\oplus\mathbb{C}E_{ji})$. A raising operator (resp. lowering operator) is an element $X\in\bigoplus_{\alpha\in\Phi_{\mathfrak{g}}^{+}}\mathfrak{g}_{\alpha}$ (resp. $X\in\bigoplus_{\alpha\in\Phi_{\mathfrak{g}}^{+}}\mathfrak{g}_{-\alpha}$). Let $\mathfrak{n}=\bigoplus_{\alpha\in\Phi_{\mathfrak{g}}^{+}}$ be the space of raising operators; the condition that a root $\alpha_{1}+\alpha_{2}$ is positive whenever $\alpha_{1},\alpha_{2}$ are positive make $\mathfrak{n}$ into a Lie algebra.





For $\mathfrak{gl}_{k}$, the space $\mathfrak{n}=\textup{span}\{E_{i,j}:i<j\}$ of raising operators is the space of strictly upper triangular matrices and the space of lowering operators is the space of strictly lower triangular matrices.





Let $V$ be a representation of $\mathfrak{g}$. A highest weight vector for $V$ is a weight vector $v\in V$ with the property that $\mathfrak{n}.v=0$. The corresponding weight is called a highest weight for $V$. When $\mathfrak{g}=\mathfrak{gl}_{k}$ and $V$ is a space of metapolynomials, it is not hard to verify that this condition is equivalent to the characterization of highest weight metapolynomials in Section 1; the characterization in terms of tableaux is explained in details in Section 4.5. A fundamental result in representation theory is that if $V$ is irreducible, then $V$ has a unique highest weight and this highest weight uniquely determines $V$ up to isomorphism [32, Prop. 14.13]. In particular, this result classifies finite-dimensional representations of $\mathfrak{g}$.





Only a subset of weights occur as highest weights of finite-dimensional irreducible representations. Such vectors are called dominant weights and their characterization in general requires the introduction of a suitable inner product structure on $\mathfrak{h}^{*}$, which is not straightforward. Let $\Lambda^{+}$ denote the subset of the weight lattice consisting of dominant weights. The representation associated to the highest weight $\lambda$ is denoted $V_{\lambda}$ and it is the unique $\mathfrak{g}$-representation having highest weight $\lambda$. In the case of $\mathfrak{gl}_{k}$, the set of dominant weights consist of weights $\lambda=(\lambda_{1},\ldots,\lambda_{k})$ such that $\lambda_{j}\geq\lambda_{j+1}$. If $\lambda_{i}\geq 0$ for every $i$, we usually identify the weight $\lambda$ with a partition of length $k$.





\par\medskip\noindent\textbf{Example 4.8.}\par



We determine the highest weights vectors in the representations $\mathbb{C}[x_{1},x_{2},x_{3}]_{2}$ and $\mathbb{C}[\mathbb{C}[x_{1},x_{2},x_{3}]_{2}]_{3}$. The space of raising operators of $\mathfrak{gl}_{3}$ is $\mathfrak{n}=\textnormal{span}\{E_{12},E_{23},E_{13}\}$ of upper triangular matrices. The highest weight vectors are weight vectors characterized by the condition $\mathfrak{n}.v=0$. Since $E_{13}=[E_{12},E_{23}]$, the condition $\mathfrak{n}.v=0$ is equivalent to $E_{12}.v=E_{23}.v=0$.





For $\mathbb{C}[x_{1},x_{2},x_{3}]_{2}$, one can check the unique highest weight vector is $x_{3}^{2}$ and the corresponding highest weight is $(0,0,-2)$. Therefore $\mathbb{C}[x_{1},x_{2},x_{3}]_{2}$ is irreducible, and it is isomorphic to the $\mathfrak{gl}_{3}$-representation $V_{(0,0,-2)}$.





For $\mathbb{C}[\mathbb{C}[x,y,z]_{2}]_{3}$, the highest weight vectors are the metapolynomials



\[\begin{aligned}
\begin{gathered}c_{2,0,0}^{3},\qquad\quad c_{2,0,0}c_{1,1,0}^{2}-4c_{2,0,0}^{2}c_{0,2,0}\qquad\quad\text{and}\\ 4c_{2,0,0}c_{0,2,0}c_{0,0,2}-c_{0,0,2}c_{1,1,0}^{2}+c_{1,0,1}c_{1,1,0}c_{0,1,1}-c_{0,2,0}c_{1,0,1}^{2}-c_{2,0,0}c_{0,1,1}^{2}.\end{gathered}
\end{aligned}\]


Their weights are $(6,0,0)$, $(4,2,0)$ and $(2,2,2)$. Therefore $\mathbb{C}[\mathbb{C}[x_{1},x_{2},x_{3}]_{2}]_{3}$ is a direct sum of three irreducible $\mathfrak{gl}_{3}$-representations



\[\begin{aligned}
\mathbb{C}[\mathbb{C}[x_{1},x_{2},x_{3}]_{2}]_{3}=V_{(6,0,0)}\ \oplus\ V_{(4,2,0)}\ \oplus\ V_{(2,2,2)}.\ 
\end{aligned}\]






\par\medskip\noindent\textbf{4.3 Universal enveloping algebra}\par




The universal enveloping algebra of a Lie algebra $\mathfrak{g}$ is an associative algebra associated to the Lie algebra $\mathfrak{g}$. It is used to study representation theory of Lie algebras using representation theory of associative algebras.





A representation of an associative (unital) algebra $\mathcal{A}$ is a vector space $V$ equipped with a homomorphism of associative algebras $\rho\colon\mathcal{A}\to\textnormal{End}(V)$. For $v\in V$ and $U\in\mathcal{A}$, write $U.v=\rho(U)v$, as in the Lie algebra setting. Subrepresentations and equivariant maps are defined for representations of an associative algebra in the same way as for Lie algebra representations. For a subset $S\subseteq\mathcal{A}$, the subalgebra generated by $S$ is the smallest subalgebra of $\mathcal{A}$ containing $S$; explicitly, it can be realized as the set of (noncommutative) polynomials in elements of $S$.





Let $\mathfrak{g}$ be a Lie algebra. Consider a representation $\rho\colon\mathfrak{g}\to\textnormal{End}(V)$.
Since $\textnormal{End}(V)$ is an associative algebra, we can consider the subalgebra of $\textnormal{End}(V)$ generated by the linear maps in the image $\rho(\mathfrak{g})$.
Since $\rho$ is a representation, these generators satisfy the relations



\[\begin{aligned}
\rho(X)\rho(Y)-\rho(Y)\rho(X)=\rho([X,Y]).\qquad\text{(6)}
\end{aligned}\]




The universal enveloping algebra $\mathcal{U}(\mathfrak{g})$ is an associative algebra constructed in such a way that the relations (6) are forced for every representation $\rho$ of $\mathcal{U}(\mathfrak{g})$.
It is a quotient of the tensor algebra $\mathcal{T}(\mathfrak{g})$ over $\mathfrak{g}$. The tensor algebra is spanned by formal products of elements of $\mathfrak{g}$,



\[\begin{aligned}
\displaystyle\mathcal{T}(\mathfrak{g})=\bigoplus_{k=0}^{\infty}\mathfrak{g}^{\otimes k}=\mathbb{C}\oplus\mathfrak{g}\oplus\mathfrak{g}^{\otimes 2}\oplus\mathfrak{g}^{\otimes 3}\oplus\cdots\qquad\text{(7)}
\end{aligned}\]


with multiplication induced by the tensor products $\otimes\colon\mathfrak{g}^{\otimes p}\times\mathfrak{g}^{\otimes q}\to\mathfrak{g}^{\otimes(p+q)}$. In particular, the tensor algebra is graded, with $\mathfrak{g}^{\otimes p}$ being the component of degree $p$. Consider the ideal $\mathcal{I}(\mathfrak{g})\subset\mathcal{T}(\mathfrak{g})$ generated by the elements of the form



\[\begin{aligned}
X\otimes Y-Y\otimes X-[X,Y]
\end{aligned}\]


for $X,Y\in\mathfrak{g}$.
The universal enveloping algebra of $\mathfrak{g}$ is defined as $\mathcal{U}(\mathfrak{g})=\mathcal{T}(\mathfrak{g})/\mathcal{I}(\mathfrak{g})$.
The multiplication in $\mathcal{U}(\mathfrak{g})$ is denoted by the usual multiplication sign rather than $\otimes$.
The inclusion of $\mathfrak{g}$ into $\mathcal{T}(\mathfrak{g})$ (as a vector space) gives rise to a linear map $\mathfrak{g}\to\mathcal{U}(\mathfrak{g})$, which is also injective.
We identify $\mathfrak{g}$ with the image of this injection.
The elements of $\mathcal{T}(\mathfrak{g})$ and, therefore, $\mathcal{U}(\mathfrak{g})$ can be written as noncommutative polynomials in elements of $\mathfrak{g}$. When $\mathfrak{g}$ is abelian (for instance, in the case of a Cartan algebra), the ideal $\mathcal{I}(\mathfrak{g})$ is indeed homogeneous, since $[X,Y]=0$; in this case $\mathcal{U}(\mathfrak{g})$ is the symmetric algebra of $\mathfrak{g}$, isomorphic to the polynomial ring on $\mathfrak{g}^{*}$.





By construction of the ideal $\mathcal{I}(\mathfrak{g})$ we have $XY-YX=[X,Y]$ in $\mathcal{U}(\mathfrak{g})$ for every $X,Y\in\mathfrak{g}$. This implies that every representation of $\mathcal{U}(\mathfrak{g})$ restricts to a Lie algebra representation of the Lie algebra $\mathfrak{g}$, via the embedding $\mathfrak{g}\subseteq\mathcal{U}(\mathfrak{g})$; conversely every representation $\rho\colon\mathfrak{g}\to\textnormal{End}(V)$ of $\mathfrak{g}$ uniquely extends to a representation of $\mathcal{U}(\mathfrak{g})$ by setting $\rho(X_{1}X_{2}\dots X_{\ell})=\rho(X_{1})\rho(X_{2})\dots\rho(X_{\ell})$ for all $X_{i}\in\mathfrak{g}$.





Note that in general the ideal $\mathcal{I}(\mathfrak{g})$ is not homogeneous, so $\mathcal{U}(\mathfrak{g})$ does not inherit the grading from the tensor algebra $\mathcal{T}(\mathfrak{g})$. However, $\mathcal{U}(\mathfrak{g})$ has a filtration, in the following sense: define $\mathcal{U}(\mathfrak{g})_{\leq\ell}=\{X+\mathcal{I}(\mathfrak{g})\mid X\in\bigoplus_{k=0}^{\ell}\mathfrak{g}^{\otimes k}\}$; then $\mathcal{U}(\mathfrak{g})_{\leq\ell}\subseteq\mathcal{U}(\mathfrak{g})_{\leq(\ell+1)}$ for every $\ell$, and if $X\in\mathcal{U}(\mathfrak{g})_{\leq p},Y\in\mathcal{U}(\mathfrak{g})_{\leq q}$ then $XY\in\mathcal{U}(\mathfrak{g})_{\leq(p+q)}$.





In particular, the elements of $\mathcal{U}(\mathfrak{g})_{\leq\ell}$ are represented by noncommutative polynomials of degree at most $\ell$ in elements of $\mathfrak{g}$.
The length of an element $X\in\mathcal{U}(\mathfrak{g})$ is the minimal $\ell$ such that $X\in\mathcal{U}(\mathfrak{g})_{\leq\ell}$.





\par\medskip\noindent\textbf{Theorem 4.9 (Poincaré–Birkhoff–Witt [42, Theorem 9.10]).}\par



Let $X_{1},\dots,X_{r}$ be a basis of $\mathfrak{g}$. Then the set of ordered monomials $X_{i_{1}}X_{i_{2}}\cdots X_{i_{d}}$, $i_{1}\leq i_{2}\leq\cdots\leq i_{d}$ is a basis of $\mathcal{U}(\mathfrak{g})$.
Moreover, the set of ordered monomials of length at most $\ell$ is a basis of the subspace $\mathcal{U}(\mathfrak{g})_{\leq\ell}$ of elements with length at most $\ell$.






In fact, we will only use that ordered monomials are a set of generators for the spaces $\mathcal{U}(\mathfrak{g})_{\leq\ell}$. This is proved explicitly in Theorem B.1.
The Poincaré–Birkhoff–Witt theorem (PBW theorem) shows that elements of $\mathcal{U}(\mathfrak{g})$ have unique representation as noncommutative polynomials in $X_{1},\dots,X_{r}$ with ordered terms, and the length of the element is given by the longest monomial in this polynomial.





\par\medskip\noindent\textbf{Example 4.10.}\par



Consider $\mathfrak{g}=\mathfrak{gl}_{2}$ with a $\mathfrak{gl}_{2}$-representation $V$.
Then $F\coloneqq E_{2,1},E\coloneqq E_{1,2},H_{i}\coloneqq E_{i,i}$ gives an ordered basis $(F,E,H_{1},H_{2})$ of $\mathfrak{gl}_{2}$.
Examples of elements of $\mathcal{U}(\mathfrak{gl}_{2})$ are $H_{1}EF$ and $1-\tfrac{1}{2}FE+\tfrac{1}{12}F^{2}E^{2}$.
Note that these are formal products, not matrix multiplication (general Lie algebras do not have this additional algebra structure).
Taking $v\in V$, the second element acts as
$(1-\tfrac{1}{2}FE+\tfrac{1}{12}F^{2}E^{2}).v=v-\tfrac{1}{2}F.(E.v)+\tfrac{1}{12}F.(F.(E.(E.v)))\in V$.
The PBW theorem guarantees that these elements can be written as polynomials with ordered monomials. Indeed, we have



\[\begin{aligned}
\displaystyle H_{1}EF & \displaystyle=H_{1}(FE+[E,F])=H_{1}FE+H_{1}(H_{1}-H_{2}) \\
 & \displaystyle=\cdots=FEH_{1}+2FE+H_{1}^{2}-H_{1}H_{2}.\ 
\end{aligned}\]






\par\medskip\noindent\textbf{4.4 Central characters and Casimir elements}\par




Let $\mathfrak{g}$ be a Lie algebra. An element $X\in\mathcal{U}(\mathfrak{g})$ of the universal enveloping algebra is central if $XY=YX$ for every $Y\in\mathcal{U}(\mathfrak{g})$ (or equivalently every $Y\in\mathfrak{g}$). The center $\mathcal{Z}(\mathfrak{g})$ of $\mathcal{U}(\mathfrak{g})$ is the set of all central elements; $\mathcal{Z}(\mathfrak{g})$ is a commutative associative subalgebra of $\mathcal{U}(\mathfrak{g})$.





A representation $\rho\colon\mathfrak{g}\to\textnormal{End}(V)$ of $\mathfrak{g}$ extends to a representation of $\mathcal{U}(\mathfrak{g})$, and if $X\in\mathcal{Z}(\mathfrak{g})$ is a central element, the map $\rho(X):V\to V$ is $\mathcal{U}(\mathfrak{g})$-equivariant. If $V$ is irreducible, then by Schur’s lemma (Lemma 4.4) $\rho(X)$ is a multiplication by some scalar $\chi_{V}(X)$. More precisely, $\rho$ induces a $1$-dimensional representation $\chi_{V}:\mathcal{Z}(\mathfrak{g})\to\mathbb{C}$. This representation is called the central character of $V$.





It is a crucial fact that for Lie algebras $\mathfrak{g}$ of reductive groups central characters separate non-isomorphic irreducible representations.





\par\medskip\noindent\textbf{Theorem 4.11 ([65], see also [10, §VIII.8.5, Cor.2]).}\par



Let $\mathfrak{g}$ be a reductive Lie algebra.
If $V$ and $W$ are irreducible representations of $\mathfrak{g}$ and the central characters $\chi_{V}$ and $\chi_{W}$ coincide, then $V$ and $W$ are isomorphic.






It follows that for reductive $\mathfrak{g}$ the isotypic components of a $\mathfrak{g}$-representation are also isotypic components of the induced $\mathcal{Z}(\mathfrak{g})$-representation.
More specifically, if a $\mathfrak{g}$-representation $V$ decomposes as $\bigoplus_{\lambda\in\Lambda}V_{\lambda}^{\oplus m_{\lambda}}$ with pairwise non-isomorphic $V_{\lambda}$, then as a $\mathcal{Z}(\mathfrak{g})$-representation we have $V\cong\bigoplus_{\lambda\in\Lambda}\chi_{\lambda}^{\oplus m_{\lambda}\cdot\dim V_{\lambda}}$, where we denote by $\chi_{\lambda}$ the $1$-dimensional representation of $\mathcal{Z}(\mathfrak{g})$ corresponding to the central character of $V_{\lambda}$.





\par\medskip\noindent\textbf{Example 4.12.}\par



Consider the $\mathfrak{gl}_{3}$ representation $V=\mathbb{C}[x_{1},x_{2},x_{3}]_{d}$ of homogeneous degree 3 polynomials in three variables (˜4.1). Consider the element $I\coloneqq-(E_{1,1}+E_{2,2}+E_{3,3})$.
Clearly, $I\in\mathcal{Z}(\mathfrak{gl}_{3})$. We determine the value of the central characters. Take $f\in V$.
We find



\[\begin{aligned}
\displaystyle I.f=x_{1}\frac{\partial}{\partial x_{1}}f+x_{2}\frac{\partial}{\partial x_{2}}f+x_{3}\frac{\partial}{\partial x_{3}}f=df.
\end{aligned}\]


So $\chi_{V}(I)=d$.
As a simple applications of Theorem 4.11, this allows us to conclude $\mathbb{C}[x_{1},x_{2},x_{3}]_{d}$ and $\mathbb{C}[x_{1},x_{2},x_{3}]_{d^{\prime}}$ contain no isomorphic irreducible subrepresentations unless $d=d^{\prime}$.






When $\mathfrak{g}$ is the Lie algebra of a reductive group, then $\mathcal{Z}(\mathfrak{g})$ can be described explicitly as the ring of invariants for the the action of the Weyl group on the Cartan subalgebra.
Note that since the Cartan subalgebra $\mathfrak{h}$ is abelian, the universal enveloping algebra $\mathcal{U}(\mathfrak{h})$ is in fact commutative and can be identified with the polynomial algebra $\mathbb{C}[\mathfrak{h}^{*}]$.
Recall that the Weyl group $\mathcal{W}_{\mathfrak{g}}$ acts on $\mathfrak{h}$ and, therefore, on $\mathcal{U}(\mathfrak{h})\cong\mathbb{C}[\mathfrak{h}^{*}]$.





\par\medskip\noindent\textbf{Theorem 4.13 (Harish–Chandra isomorphism, [44, Ch. 23];[6]).}\par



Let $\mathfrak{g}$ be a reductive algebra and let $\mathfrak{h}$ be a Cartan subalgebra.
The center $\mathcal{Z}(\mathfrak{g})$ of the universal enveloping algebra is isomorphic to the algebra of invariants $\mathcal{U}(\mathfrak{h})^{\mathcal{W}_{\mathfrak{g}}}$. Moreover, if $\dim\mathfrak{h}=k$, there exist algebraically independent $C_{1},\dots,C_{k}\in\mathcal{U}(\mathfrak{g})$ such that $\mathcal{Z}(\mathfrak{g})=\mathbb{C}[C_{1},\dots,C_{k}]$. The lengths of these elements coincide with the degrees of the fundamental invariants of $\mathcal{U}(\mathfrak{h})^{\mathcal{W}_{\mathfrak{g}}}$.






The generators $C_{1},\dots,C_{k}$ of the center $\mathcal{Z}(\mathfrak{g})$ are known as Casimir elements.
These generators are not unique.





Consider the case of $\mathfrak{gl}_{k}$, with $\mathfrak{h}$ the Cartan subalgebra of diagonal matrices. Then $\mathcal{W}_{\mathfrak{gl}_{k}}$ is the symmetric group $\mathfrak{S}_{k}$, and it acts on $\mathfrak{h}$ by permuting the diagonal elements. Then $\mathcal{U}(\mathfrak{h})$ is the polynomial ring in $k$ variables, and $\mathcal{U}(\mathfrak{h})^{\mathcal{W}_{\mathfrak{gl}_{k}}}$ is the subring of polynomials invariant under permutation of variables: this is the classical ring of symmetric polynomials in $k$ variables. It is generated, for instance, by the elementary symmetric polynomials, which have degrees $1,\ldots,k$. These numbers give lengths of Casimir elements generating the center $\mathcal{Z}(\mathfrak{g})$.
There are several constructions of Casimir elements, one possible choice being



\[\begin{aligned}
\displaystyle C_{p}\coloneqq\sum_{i_{1}=1}^{k}\dots\sum_{i_{p}=1}^{k}E_{i_{1},i_{2}}E_{i_{2},i_{3}}\cdots E_{i_{p},i_{1}}\quad\text{with }1\leq p\leq k.\qquad\text{(8)}
\end{aligned}\]




\par\medskip\noindent\textbf{Example 4.14.}\par



Let $\mathfrak{g}=\mathfrak{gl}_{2}$. Then the Casimir elements are given by



\[\begin{aligned}
\displaystyle C_{1}=E_{1,1}+E_{2,2}\qquad\text{and}\qquad C_{2}=E_{1,1}E_{1,1}+E_{1,2}E_{2,1}+E_{2,1}E_{1,2}+E_{2,2}E_{2,2}.
\end{aligned}\]


We will use ˜4.2 to evaluate the application of $C_{1}$ and $C_{2}$ to the metapolynomials



\[\begin{aligned}
\displaystyle\Delta=c_{11}^{2}+2\,c_{02}c_{20}\qquad\text{and}\qquad\Gamma=-c_{11}^{2}+4\,c_{02}c_{20}.
\end{aligned}\]


We find that



\[\begin{aligned}
\displaystyle C_{1}.\Delta=E_{1,1}.\Delta+E_{2,2}.\Delta=2c_{11}^{2}+4c_{02}c_{20}+2c_{11}^{2}+4c_{02}c_{20}=4\Delta
\end{aligned}\]


and similarly $C_{1}.\Gamma=4\Gamma$ (in fact, one can show $C_{1}$ always scales a metapolynomial of format $(\delta,d,k)$ by $d\delta$).
We also compute



\[\begin{aligned}
\displaystyle E_{2,1}.\Delta & \displaystyle=2\,c_{02}^{2}c_{11}+2\,c_{11}^{2}c_{02}+2\cdot 0c_{20}+2\,c_{02}c_{11}=6\,c_{02}^{2}c_{11} \\
\displaystyle(E_{1,2}E_{2,1}).\Delta & \displaystyle=6\,c_{11}c_{11}+6\cdot 2\,c_{02}c_{20}=6\Delta
\end{aligned}\]


By symmetry of multiindices in the metapolynomial we know $(E_{2,1}E_{1,2}).\Delta=6\Delta$ too.
An analogous computation shows $E_{2,1}.\Gamma=0$ and $E_{1,2}.\Gamma=0$.
We find



\[\begin{aligned}
\displaystyle C_{2}.\Delta=\big{(}2^{2}+6+6+2^{2}\big{)}\Delta=20\Delta\qquad\text{and}\qquad C_{2}.\Gamma=\big{(}2^{2}+0+0+2^{2}\big{)}\Delta=8\Delta.
\end{aligned}\]


We obtain that $\Delta$ and $\Gamma$ live in different isotypic components with respect to
$\mathcal{Z}(\mathfrak{gl}_{2})$, and hence with respect to $\mathfrak{gl}_{2}$.






The eigenvalues of these Casimir elements on irreducible representations of $\mathfrak{gl}_{k}$ are computed by Perelomov and Popov [64]:



\[\begin{aligned}
\displaystyle\chi_{\lambda}(C_{p})=\mathsf{Tr}(A_{\lambda}^{p}E)\qquad\text{(9)}
\end{aligned}\]


where $E$ is the $k\times k$ matrix consisting of all ones, and the entries of the $k\times k$ matrix $A_{\lambda}=(a^{(\lambda)}_{ij})$ are given by



\[\begin{aligned}
\displaystyle a^{(\lambda)}_{ij}=\begin{cases}\lambda_{i}+k-i,&\text{if $i=j$,}\\ 0,&\text{if $i>j$,}\\ -1,&\text{if $i<j$.}\end{cases}\qquad\text{(10)}
\end{aligned}\]




\par\medskip\noindent\textbf{Example 4.15.}\par



Observe first that $\mathsf{Tr}(A_{\lambda}^{p}E)$ is simply the sum of the entries of $A_{\lambda}^{p}$.
Consider again Example 4.14. We compute the formula in Equation 9.



\[\begin{aligned}
\displaystyle A_{(4,0)} & \displaystyle=\left[\begin{smallmatrix}5&-1\\ 0&2\\ \end{smallmatrix}\right]\ \text{ and thus }\ \chi_{(4,0)}(C_{2})=\mathsf{Tr}\left(\left[\begin{smallmatrix}25&-7\\ 0&4\\ \end{smallmatrix}\right]\left[\begin{smallmatrix}1&1\\ 1&1\\ \end{smallmatrix}\right]\right)=20 \\
\displaystyle A_{(2,2)} & \displaystyle=\left[\begin{smallmatrix}3&-1\\ 0&2\\ \end{smallmatrix}\right]\ \text{ and thus }\ \chi_{(2,2)}(C_{2})=\mathsf{Tr}\left(\left[\begin{smallmatrix}9&-5\\ 0&4\\ \end{smallmatrix}\right]\left[\begin{smallmatrix}1&1\\ 1&1\\ \end{smallmatrix}\right]\right)=8. \\
\mathfrak{gl}_{3} \\
\displaystyle A_{(6,0,0)} & \displaystyle=\left[\begin{smallmatrix}8&-1&-1\\ 0&1&-1\\ 0&0&0\end{smallmatrix}\right]\ \text{ and thus }\ \chi_{(6,0,0)}(C_{2})=\mathsf{Tr}\left(\left[\begin{smallmatrix}64&-9&-7\\ 0&1&-1\\ 0&0&0\end{smallmatrix}\right]\left[\begin{smallmatrix}1&1&1\\ 1&1&1\\ 1&1&1\end{smallmatrix}\right]\right)=48.\ 
\end{aligned}\]






\par\medskip\noindent\textbf{4.5 Gelfand–Tsetlin theory}\par




Gelfand–Tsetlin bases are used to explicitly construct irreducible representations of some classical Lie algebras.
The explicit bases first appeared in [35] without much explanation.
The construction was expanded upon by Želobenko [72] and independently by Baird and Biedenharn [4].
It can be explained in terms of restrictions of irreducible representations along a chain of Lie algebra inclusions.





Let $\mathfrak{a}$ be a Lie subalgebra of $\mathfrak{b}$.
Every representation $V$ of $\mathfrak{b}$ can be regarded as a representation of $\mathfrak{a}$ using this inclusion.
The resulting representation is called the restriction of $V$ to $\mathfrak{a}$ and is denoted as $V{\downarrow}^{\mathfrak{b}}_{\mathfrak{a}}$.





Let $\mathfrak{g}_{1}\subset\mathfrak{g}_{2}\subset\dots\subset\mathfrak{g}_{k}$ be an inclusion chain of reductive Lie algebras.
We call this chain a Gelfand–Tsetlin chain if $\mathfrak{g}_{1}$ is abelian and every irreducible representation $V$ of $\mathfrak{g}_{\ell}$ ($1<\ell\leq k$), when restricted to $\mathfrak{g}_{\ell-1}$, is a multiplicity-free completely reducible representation, that is,



\[\begin{aligned}
\displaystyle V_{\lambda}{\downarrow}^{\mathfrak{g}_{\ell}}_{\mathfrak{g}_{\ell-1}}\cong\bigoplus_{\mu\in R_{\ell}(\lambda)}V_{\mu}\qquad\text{(11)}
\end{aligned}\]


for some set $R_{\ell}(\lambda)$ of isomorphism types of irreducible representations of $\mathfrak{g}_{i-1}$.
We write $\mu\to\lambda$ if $V_{\mu}$ appears in the decomposition for $V_{\lambda}$.





Let $V_{\lambda}$ be a irreducible representation of $\mathfrak{g}_{k}$.
By repeatedly restricting irreducible representations of $\mathfrak{g}_{i}$ we obtain a decomposition of $V_{\lambda}$ as a representation of $\mathfrak{g}_{1}$.



\[\begin{aligned}
\displaystyle V_{\lambda}{\downarrow}^{\mathfrak{g}_{k}}_{\mathfrak{g}_{1}}=\bigoplus_{\lambda^{(1)}\to\lambda^{(2)}\to\dots\to\lambda^{(k)}=\lambda}V_{\lambda^{(1)}}\qquad\text{(12)}
\end{aligned}\]


Since $\mathfrak{g}_{1}$ is abelian, its irreducible representations are $1$-dimensional, so the decomposition above distinguishes a set of $1$-dimensional subspaces in $V_{\lambda}$ indexed by chains $\lambda^{(1)}\to\lambda^{(2)}\to\dots\to\lambda^{(k)}$ called Gelfand–Tsetlin patterns.
We denote the subspace of $V_{\lambda^{(k)}}$ corresponding to a Gelfand–Tsetlin pattern $T=\lambda^{(1)}\to\lambda^{(2)}\to\dots\to\lambda^{(k)}$ by $V_{T}$.





The inclusions $\mathfrak{g}_{\ell}\subseteq\mathfrak{g}_{k}$ give an inclusion of universal enveloping algebras $\mathcal{U}(\mathfrak{g}_{\ell})\subseteq\mathcal{U}(\mathfrak{g}_{k})$, and therefore a chain $\mathcal{U}(\mathfrak{g}_{1})\subset\mathcal{U}(\mathfrak{g}_{2})\subset\dots\subset\mathcal{U}(\mathfrak{g}_{k})$. The Gelfand–Tsetlin algebra $\mathcal{G}\hskip-1.00006pt\mathcal{Z}(\mathfrak{g}_{k})$ of $\mathfrak{g}_{k}$ is the subalgebra of $\mathcal{U}(\mathfrak{g})$ generated by all centers $\mathcal{Z}(\mathfrak{g}_{\ell})$ with $\ell\leq k$. Since the center $\mathcal{Z}(\mathfrak{g}_{\ell})$ of $\mathcal{U}(\mathfrak{g}_{\ell})$ commutes with all elements of $\mathcal{U}(\mathfrak{g}_{\ell})$, it commutes in particular with the centers $\mathcal{Z}(\mathfrak{g}_{j})$ for $j<\ell$; as a result, $\mathcal{G}\hskip-1.00006pt\mathcal{Z}(\mathfrak{g}_{k})$ is a commutative subalgebra.





Consider a Gelfand–Tsetlin pattern $T=\lambda^{(1)}\to\lambda^{(2)}\to\dots\to\lambda^{(k)}$ and the corresponding $1$-dimensional subspace $V_{T}\subset V_{\lambda^{(k)}}$.
For every $\ell\leq k$ the subspace $V_{T}$ is contained in an irreducible representation of $\mathfrak{g}_{\ell}$ isomorphic to $V_{\lambda^{(\ell)}}$, and the center $\mathcal{Z}(\mathfrak{g}_{\ell})$ acts on $V_{T}$ via the central character $\chi_{\lambda^{(\ell)}}$.
It follows that $V_{T}$ is a $1$-dimensional representation of $\mathcal{G}\hskip-1.00006pt\mathcal{Z}(\mathfrak{g}_{k})$ given by a Gelfand–Tsetlin character $\chi_{T}\colon\mathcal{G}\hskip-1.00006pt\mathcal{Z}(\mathfrak{g}_{k})\to\mathbb{C}$ such that $X.v=\chi_{T}(X)v$ for $X\in\mathcal{G}\hskip-1.00006pt\mathcal{Z}(\mathfrak{g}_{k})$, $v\in V_{T}$.
More specifically, on $\mathcal{Z}(\mathfrak{g}_{\ell})$ we have $\chi_{T}(X)=\chi_{\lambda^{(\ell)}}(X)$.
Theorem 4.11 implies that different Gelfand–Tsetlin patterns correspond to different Gelfand–Tsetlin characters.





By the previous discussion, every completely reducible representation of $\mathfrak{g}_{k}$ is also completely reducible as a representation of $\mathcal{G}\hskip-1.00006pt\mathcal{Z}(\mathfrak{g}_{k})$.
That is, if $V$ is a completely reducible representation of $\mathfrak{g}_{k}$ with isotypic decomposition $V=\bigoplus_{\lambda\in\Lambda}V_{\lambda}^{\oplus m_{\lambda}}$, then the action of $\mathcal{G}\hskip-1.00006pt\mathcal{Z}(\mathfrak{g}_{k})$ on $V$ splits the decomposition further into Gelfand–Tsetlin characters as



\[\begin{aligned}
\displaystyle V=\bigoplus\limits_{T\colon\lambda^{(1)}\to\dots\to\lambda^{(k)},\lambda^{(k)}\in\Lambda}\chi_{T}^{\oplus m_{\smash{\lambda^{(k)}}}}.\qquad\text{(13)}
\end{aligned}\]


We call the components of this decomposition GZ-isotypic spaces of $V$.





The classical case – used to construct a basis of irreducible representations of $\mathfrak{gl}_{n}$ – corresponds to the chain of inclusions $\mathfrak{gl}_{1}\subset\mathfrak{gl}_{2}\subset\dots\subset\mathfrak{gl}_{k}$ where $\mathfrak{gl}_{k-1}$ is included in $\mathfrak{gl}_{k}$ as the subalgebra of $\mathfrak{gl}_{k}$ consisting of matrices having the $k$-th row and the $k$-th column identically equal to $0$. The Gelfand–Tsetlin algebra in this case is the polynomial algebra $\mathcal{G}\hskip-1.00006pt\mathcal{Z}(\mathfrak{g}_{k})=\mathbb{C}[C_{\ell,p}]$ generated by the Casimir elements



\[\begin{aligned}
\displaystyle C_{\ell,p}\coloneqq\sum_{i_{1}=1}^{\ell}\dots\sum_{i_{p}=1}^{\ell}E_{i_{1},i_{2}}E_{i_{2},i_{3}}\cdots E_{i_{p},i_{1}}\quad\text{with }1\leq\ell\leq k\text{ and }1\leq p\leq\ell.\qquad\text{(14)}
\end{aligned}\]




The irreducible representations of $\mathfrak{gl}_{k}$ correspond highest weights $\lambda=(\lambda_{1},\dots,\lambda_{k})$ with integer $\lambda_{1}\geq\dots\geq\lambda_{k}$.
The restriction $V_{\lambda}{\downarrow}^{\mathfrak{gl}_{k}}_{\mathfrak{gl}_{k-1}}$ is given by the classical branching rule (attributed to Schur in [58]) which states that $\mu\to\lambda$ if and only if $\mu$ and $\lambda$ interleave, that is $\lambda_{\ell}\geq\mu_{\ell}\geq\lambda_{\ell+1}$ for all $\ell$, $1\leq\ell\leq k-1$.
Moreover, $E_{k,k}$ acts on the copy of $V_{\mu}$ in $V_{\lambda}$ as multiplication by $|\lambda|-|\mu|=\sum_{\ell=1}^{k-1}(\lambda_{k}-\mu_{k})$.





A Gelfand-Tsetlin pattern is given by a sequence $\lambda^{(1)}\to\lambda^{(2)}\to\dots\to\lambda^{(k)}$, that is, by integers $\lambda^{(\ell)}_{i}$ with $1\leq\ell\leq k$, $1\leq i\leq\ell$ such that $\lambda^{(\ell)}_{i}\geq\lambda^{(\ell-1)}_{i}\geq\lambda^{(\ell)}_{i+1}$, often arranged into a triangular shape, see Figure 2.




Figure 2: Gelfand–Tsetlin pattern of $\mathfrak{gl}_{k}$.



The $1$-dimensional subspace $V_{T}$ corresponding to the pattern $T=\lambda^{(1)}\to\lambda^{(2)}\to\dots\to\lambda^{(k)}$ is spanned by a weight vector of weight $(|\lambda^{(1)}|,|\lambda^{(2)}|-|\lambda^{(1)}|,\dots,|\lambda^{(k)}|-|\lambda^{(k-1)}|)$.
The generators $C_{\ell,p}$ of $\mathcal{G}\hskip-1.00006pt\mathcal{Z}(\mathfrak{gl}_{k})$ act on this space as multiplications by the corresponding eigenvalues $\chi_{T}(C_{\ell,p})$.
Since $C_{\ell,p}$ is a Casimir operator for $\mathfrak{gl}_{\ell}$, we have $\chi_{T}(C_{\ell,p})=\chi_{\lambda^{(\ell)}}(C_{\ell,p})=\mathsf{Tr}(A_{\lambda^{(\ell)}}^{p}E)$ from Perelomov–Popov formula.





If $\lambda^{(k)}$ is a partition, that is, all $\lambda^{(k)}_{i}$ are nonnegative, then the Gelfand–Tsetlin patterns for $\mathfrak{gl}_{k}$ are in bijective correspondence with semistandard Young tableaux filled with numbers from $\{1,2,\dots,k\}$. Recall that a Young tableau is semistandard whenever the numbers are strictly increasing along the columns, and weakly increasing along the rows.





A semistandard tableau $T$ gives rise to a sequence of partitions $\lambda^{(1)},\dots,\lambda^{(k)}$ where $\lambda^{(\ell)}$ is the shape of the tableau $T^{(\ell)}$ obtained from $T$ be removing all boxes labeled with numbers greater than $\ell$.
Semistandardness of $T$ guarantees that all $T^{(\ell)}$ are also semistandard tableaux and the interleaving property $\lambda^{(\ell)}_{i}\geq\lambda^{(\ell-1)}_{i}\geq\lambda^{(\ell)}_{i+1}$ holds (see e.\thinspace{}g. [56, p.5]).
The weight is given by the content of the tableau $T$, that is, the vector $\mu=(\mu_{1},\dots,\mu_{k})$ where $\mu_{i}$ records the number of occurrences of the label $i$ in the tableau.





For example, the tableau 
gives








and the sequence of shapes is $\lambda^{(1)}=(2)$, $\lambda^{(2)}=(4,1)$, $\lambda^{(3)}=(4,2,0)$, $\lambda^{(4)}=(4,2,2,0)$.
The corresponding weight is $(2,3,1,2)$, the content of the tableau.





The superstandard tableau of shape $\lambda=(\lambda_{1},\dots,\lambda_{n})$ corresponds to the Gelfand-Tsetlin pattern with $\lambda^{(\ell)}=(\lambda_{1},\dots,\lambda_{\ell})$.
The corresponding $1$-dimensional subspace of $V_{\lambda}$ contains the highest weight vectors.
To prove this, recall that since $E_{i,j}$ with $i\neq j$ span root spaces of $\mathfrak{gl}_{k}$, they shift the weights of weight vectors, that is, $E_{i,j}.v$ is a weight vector of weight $\mu+e_{i}-e_{j}$.
The content of the superstandard tableau of shape $\lambda$ is $\lambda$, so every nonzero vector $v$ in the corresponding $1$-dimensional subspace has weight $\lambda$, and for every $E_{i,j}$ the vector $E_{i,j}.v$ has weight $\lambda+e_{i}-e_{j}$. But there are no semistandard tableaux with shape $\lambda$ and content $\lambda+e_{i}-e_{j}$ when $i<j$.
Therefore, $E_{i,j}.v=0$ if $i<j$, and it follows that $v$ is annihilated by every strictly upper triangular matrix.




